\documentclass[11pt,a4paper]{article}

% Même préambule que la proposition sous XeLaTeX/LuaLaTeX ; repli pdflatex
% (mêmes polices Latin Modern) quand fontspec n'est pas utilisable.
\usepackage{iftex}
\ifPDFTeX
  \usepackage[T1]{fontenc}
  \usepackage[utf8]{inputenc}
  \usepackage{lmodern}
  \usepackage[french]{babel}
  \frenchsetup{StandardLists=true, AutoSpacePunctuation=false, og=«, fg=»}
\else
  \usepackage{fontspec}
  \usepackage{polyglossia}
  \setmainlanguage{french}
\fi
\usepackage[a4paper,margin=2.5cm]{geometry}
\usepackage{microtype}
\usepackage{amsmath,amssymb}
\usepackage{booktabs,tabularx,array}
\usepackage{enumitem}
\usepackage{titlesec}
\usepackage{fancyhdr}
\usepackage{xcolor}
\usepackage{caption}
\usepackage{float}
\usepackage{graphicx}
\usepackage{tikz}
\usetikzlibrary{positioning,arrows.meta,fit,backgrounds,shapes.geometric,calc}
\usepackage[hidelinks]{hyperref}
\renewcommand{\topfraction}{0.9}
\renewcommand{\bottomfraction}{0.7}
\renewcommand{\textfraction}{0.1}
\renewcommand{\floatpagefraction}{0.75}

\setlength{\parskip}{0.5em}
\setlength{\parindent}{0pt}
\setlist{itemsep=0.15em,topsep=0.3em}
\captionsetup{font=small,labelfont=bf}
\renewcommand{\arraystretch}{1.25}
\newcolumntype{L}{>{\raggedright\arraybackslash}X}
\newcolumntype{P}[1]{>{\raggedright\arraybackslash}p{#1}}
\setlength{\emergencystretch}{3em}

\titleformat{\section}{\large\bfseries}{\thesection.}{0.6em}{}
\titleformat{\subsection}{\normalsize\bfseries}{\thesubsection}{0.6em}{}

\pagestyle{fancy}
\fancyhf{}
\lhead{\small INPT -- Smart ICT}
\rhead{\small IA agentique pour le RIC en O-RAN}
\cfoot{\small\thepage}
\renewcommand{\headrulewidth}{0.4pt}

\newcommand{\tool}[1]{\texttt{#1}}

\tikzset{
  box/.style={draw, rounded corners=2pt, align=left, inner sep=5pt, font=\small},
  cbox/.style={draw, rounded corners=2pt, align=center, inner sep=5pt, font=\small},
  arr/.style={-{Stealth[length=2.2mm]}, thick},
  lbl/.style={font=\footnotesize\itshape, inner sep=2pt},
  zone/.style={draw, dashed, rounded corners=3pt, fill=gray!6},
  zlbl/.style={font=\footnotesize\itshape, anchor=north west},
  off/.style={draw=gray!55, text=gray!60},
}

\begin{document}

%----------------------------------------------------------------- en-tête
\thispagestyle{plain}
\begin{center}
  {\small Institut National des Postes et Télécommunications (INPT), Rabat\\
  Filière Smart ICT -- 3\textsuperscript{e} année}\\[1.2em]
  {\LARGE\bfseries IA agentique pour le RIC en O-RAN}\\[0.4em]
  {\large Rapport d'avancement -- phases 1 à 3}\\[1.2em]
  \begin{tabular}{@{}l@{\hspace{3em}}l@{}}
    \textit{Réalisé par :} & \textit{Encadré par :}\\
    Souhail Bourhim & Pr. Mustapha Benjillali\\
    Ghita Boulmani & \\
  \end{tabular}\\[1em]
  {\small Novembre 2026}
\end{center}
\vspace{0.5em}
\hrule
\vspace{0.8em}

\textbf{Résumé.} Ce rapport présente l'avancement du projet par rapport à la proposition
d'implémentation d'octobre 2026. Il couvre les trois premières phases du plan de travail :
les fondations (simulateur, plateforme OAI + FlexRIC, supervision KPM), le contrôle (politique
RL, contrôle des PRB par E2) et l'agent (service LangGraph, vérificateur déterministe, console
de confirmation). L'interface opérateur (Hermes Agent et la passerelle à trois outils) et
l'évaluation comparative des contrôleurs feront l'objet du prochain rapport.

En résumé, la chaîne « intention $\to$ service LangGraph $\to$ vérificateur $\to$ politique A1
$\to$ xApp RL $\to$ répartition des PRB » fonctionne de bout en bout, sur le simulateur et sur une
vraie cellule OAI pilotée par E2. La politique RL est conditionnée par l'objectif et passe un test
conçu pour pouvoir échouer. Pour y arriver sur la plateforme, nous avons dû corriger le gNB OAI
à deux endroits : un bug des mesures KPM par slice, et l'absence de contrôle des PRB par E2 dans
la version utilisée. Nous indiquons aussi ce qui ne marche pas encore, en particulier une
oscillation de la politique sur la plateforme réelle (section~\ref{sec:limites}).

%----------------------------------------------------------------- 1
\section{Rappel de l'objectif et état d'avancement}

La question du projet reste la même : \emph{comment permettre à un opérateur de piloter le RAN
en langage naturel, tout en gardant des décisions rapides et en empêchant l'IA de prendre des
actions dangereuses ?} La réponse proposée sépare les rôles : un LLM traduit l'intention en
objectif, un vérificateur déterministe contrôle cet objectif, et une politique RL décide seule
de la répartition des PRB, sans jamais attendre le LLM.

La figure~\ref{fig:plan} reprend le plan de travail de la proposition. Les phases 1 à 3 sont
terminées, avec leurs points de validation ; le dernier (« intention satisfaite de bout en bout,
sans Hermes ») a été atteint sur la plateforme réelle (section~\ref{sec:boucle}).

Le travail a été découpé en étapes courtes, chacune avec un critère de réussite fixé avant de
commencer, des tests automatiques et une note technique. Le tableau~\ref{tab:etapes} liste les
étapes couvertes ici.

\begin{table}[!htb]
\centering\small
\begin{tabularx}{\textwidth}{@{}P{1.9cm}LP{1.1cm}@{}}
\toprule
Étape & Contenu & Phase\\
\midrule
M01 & Simulateur gym d'une cellule à trois slices ; schéma de l'objectif RL & 1\\
M06a & Première mise en route d'OAI + FlexRIC, capacité mesurée & 1\\
M06 & Trois slices de bout en bout sur la plateforme (c\oe ur, gNB, un UE par slice) & 1\\
M07a & xApp de supervision KPM ; correctif d'un bug du gNB & 1\\
M02, M02b & Politique RL conditionnée par l'objectif, test fondé sur un oracle & 2\\
M02c & Suivi des entraînements (MLflow), registre de modèles, intégration continue & 2\\
M07b & Contrôle des PRB par E2 ; xApp RL sur la plateforme réelle & 2, 3\\
M03 & Vérificateur déterministe & 3\\
M04a & Intégration : \tool{Backend}, base partagée, politique A1, xApp RL, console & 3\\
M04b & Service LangGraph avec un LLM local & 3\\
\bottomrule
\end{tabularx}
\caption{Étapes réalisées (M06a désigne la courte étude préalable de la plateforme).}
\label{tab:etapes}
\end{table}

\begin{figure}[H]
\centering
\begin{tikzpicture}[x=1cm,y=1cm]
  \node[font=\small\bfseries] at (6.0,0.5) {Axe A : plateforme, xApps};
  \node[font=\small\bfseries] at (12.2,0.5) {Axe B : RL, agent, interface};

  \node[align=left, anchor=west, font=\small] at (0,-0.7) {\textbf{Phase 1}\\\footnotesize Fondations};
  \node[cbox, fill=gray!20, text width=5.5cm, minimum height=1.1cm] at (6.0,-0.7) {\footnotesize Installer OAI, c\oe ur, FlexRIC ;\\\footnotesize xApp de supervision (KPM)};
  \node[cbox, fill=gray!20, text width=5.5cm, minimum height=1.1cm] at (12.2,-0.7) {\footnotesize Simulateur gym des 3 slices};
  \node[font=\footnotesize\itshape, anchor=west] at (3.1,-1.575) {$\blacklozenge$ KPI en direct \quad\textbf{atteint}};

  \node[align=left, anchor=west, font=\small] at (0,-2.45) {\textbf{Phase 2}\\\footnotesize Contrôle};
  \node[cbox, fill=gray!20, text width=5.5cm, minimum height=1.1cm] at (6.0,-2.45) {\footnotesize xApp de contrôle des slices};
  \node[cbox, fill=gray!20, text width=5.5cm, minimum height=1.1cm] at (12.2,-2.45) {\footnotesize Entraînement RL\\\footnotesize en simulation};
  \node[font=\footnotesize\itshape, anchor=west] at (3.1,-3.325) {$\blacklozenge$ Le RL modifie la répartition des PRB \quad\textbf{atteint}};

  \node[align=left, anchor=west, font=\small] at (0,-4.2) {\textbf{Phase 3}\\\footnotesize Agent};
  \node[cbox, fill=gray!20, text width=5.5cm, minimum height=1.1cm] at (6.0,-4.2) {\footnotesize Relier l'agent au RIC\\\footnotesize (politique A1, données KPM)};
  \node[cbox, fill=gray!20, text width=5.5cm, minimum height=1.1cm] at (12.2,-4.2) {\footnotesize Service LangGraph, vérificateur,\\\footnotesize console de confirmation};
  \node[font=\footnotesize\itshape, anchor=west] at (3.1,-5.075) {$\blacklozenge$ Intention satisfaite de bout en bout (sans Hermes) \quad\textbf{atteint}};

  \node[align=left, anchor=west, font=\small, text=gray!70] at (0,-5.95) {\textbf{Phase 4}\\\footnotesize Interface};
  \node[cbox, off, text width=5.5cm, minimum height=1.1cm] at (6.0,-5.95) {\footnotesize Passerelle MCP à 3 outils ;\\\footnotesize tests de refus (hors périmètre)};
  \node[cbox, off, text width=5.5cm, minimum height=1.1cm] at (12.2,-5.95) {\footnotesize Configurer Hermes ; intentions\\\footnotesize annotées ; mesures d'interface};

  \node[align=left, anchor=west, font=\small, text=gray!70] at (0,-7.7) {\textbf{Phase 5}\\\footnotesize Résultats};
  \node[cbox, off, text width=11.6cm, minimum height=1.1cm] at (9.1,-7.7) {\footnotesize Évaluation du contrôle (4 contrôleurs) et de l'interface ;\\\footnotesize analyse, rapport, préparation de la soutenance};
  \node[font=\footnotesize\itshape, text=gray!70, anchor=west] at (3.1,-6.825) {phases 4 et 5 : présentées dans le prochain rapport};
\end{tikzpicture}
\caption{Plan de travail de la proposition. En gris foncé, les phases couvertes par ce rapport.}
\label{fig:plan}
\end{figure}


%----------------------------------------------------------------- 2
\section{Architecture réalisée}

La figure~\ref{fig:archi} reprend l'architecture de la proposition. Tous les blocs du RIC et du
RAN existent. La zone de l'interface opérateur est en gris : elle relève de la phase 4. Pendant
les phases 1 à 3, les intentions sont envoyées directement à l'API HTTP du service LangGraph.

\begin{figure}[!htb]
\centering
\begin{tikzpicture}[x=1cm,y=1cm]
  % zones
  \draw[zone, draw=gray!50, fill=gray!3] (0,1.35) rectangle (15.45,-0.95);
  \node[zlbl, text=gray!60] at (0.1,1.3) {Interface opérateur (phase 4, prochain rapport)};
  \draw[zone] (5.75,-1.95) rectangle (15.4,-4.4);
  \node[zlbl] at (5.85,-2.0) {Non-RT RIC / SMO (émulé)};
  \draw[zone] (5.75,-5.15) rectangle (15.4,-7.45);
  \node[zlbl] at (8.1,-5.2) {Near-RT RIC (FlexRIC)};

  % ligne 0 : interface (hors périmètre)
  \node[box, off, text width=2.0cm] (op) at (1.35,0) {\textbf{Opérateur}\\\footnotesize langage\\\footnotesize naturel};
  \node[box, off, text width=3.0cm] (he) at (5.9,0) {\textbf{Hermes Agent}\\\footnotesize interface de chat,\\\footnotesize client MCP};
  \node[box, off, text width=4.6cm] (gw) at (12.7,0) {\textbf{Passerelle MCP}\\\footnotesize\tool{submit\_intent}, \tool{get\_kpis},\\\footnotesize\tool{get\_status}};

  % ligne 1 : Non-RT RIC
  \node[box, text width=4.0cm, fill=gray!12] (lg) at (12.9,-3.3) {\textbf{Service LangGraph}\\\footnotesize FastAPI + worker,\\\footnotesize LLM local (\tool{qwen3:8b})};
  \node[box, text width=3.2cm] (vf) at (7.7,-3.3) {\textbf{Vérificateur}\\\footnotesize Python déterministe,\\\footnotesize 8 contrôles};
  \node[box, dashed, text width=2.9cm] (cf) at (2.6,-3.3) {\textbf{Console de confirmation}\\\footnotesize démo, délai 120 s};

  % ligne 2 : Near-RT RIC
  \node[box, text width=3.4cm] (rl) at (7.7,-6.5) {\textbf{xApp RL (TD3)}\\\footnotesize décide seule de la\\\footnotesize répartition des PRB};
  \node[box, text width=4.0cm] (kp) at (12.9,-6.5) {\textbf{xApp KPM}\\\footnotesize débit, PRB, délai RLC\\\footnotesize par slice, chaque seconde};
  \node[box, text width=4.3cm] (db) at (2.65,-6.5) {\textbf{Base partagée} \tool{var/ric.db}\\\footnotesize KPI, \tool{a1\_policy} (A1 émulé),\\\footnotesize intentions ; journal JSONL};

  % RAN
  \node[box, text width=9.0cm] (ran) at (10.55,-9.3) {\textbf{RAN, derrière une seule interface \tool{Backend}}\\\footnotesize simulateur Python (\tool{SimBackend}) ou OAI : gNB, 3 nrUE, CN5G (\tool{FlexRICBackend})};

  % flèches
  \draw[arr, gray!50] (op) -- (he);
  \draw[arr, gray!50] (he) -- (gw);
  \draw[arr, gray!50] (11.5,0 |- gw.south) -- (11.5,0 |- lg.north);
  \draw[arr] (lg.west) -- node[lbl,above]{objectif} (vf.east);
  \draw[arr, <->] (vf.west) -- node[lbl,above]{démo} (cf.east);
  \draw[arr] (vf.south) -- node[lbl,right]{politique A1} (rl.north);
  \draw[arr] (kp.north) -- node[lbl,right]{KPI} (lg.south);
  \draw[arr] (kp.west) -- (rl.east);
  \draw[arr] (rl.south) -- node[lbl,right]{contrôle E2} (rl.south |- ran.north);
  \draw[arr] (kp.south |- ran.north) -- node[lbl,right]{rapports E2} (kp.south);
  \node[lbl, align=center] at (14.4,-1.45) {\tool{POST /intents}\\(phases 1--3)};
\end{tikzpicture}
\caption{Architecture réalisée. En gris : l'interface opérateur, hors du périmètre de ce rapport.
Les processus ne communiquent que par la base partagée ; la xApp RL est le seul appelant de
\tool{apply\_prb\_split}.}
\label{fig:archi}
\end{figure}

Trois choix structurent le code :
\begin{itemize}
  \item \textbf{Une seule interface pour deux RAN.} La xApp RL appelle les mêmes cinq méthodes
        (lire les KPI, appliquer une répartition, etc.) sur le simulateur et sur la plateforme.
        Le service LangGraph, le vérificateur et la politique n'ont pas changé d'une ligne pour
        passer à OAI.
  \item \textbf{Une politique A1 émulée.} FlexRIC est un Near-RT RIC : il n'y a ni Non-RT RIC
        ni interface A1 dans notre pile. L'objectif validé est donc écrit dans une table
        \tool{a1\_policy} sous forme d'un calendrier (segments de temps), que la xApp RL lit à
        chaque seconde. C'est une émulation, et nous la nommons ainsi.
  \item \textbf{Des invariants vérifiés par des tests.} Par exemple : le vérificateur ne fait
        ni appel réseau, ni tirage aléatoire, ni appel à un LLM ; la xApp RL n'importe rien du
        service LangGraph et n'attend jamais un LLM ; seule la xApp RL fixe la répartition des
        PRB.
\end{itemize}

%----------------------------------------------------------------- 3
\section{Simulateur et politique RL conditionnée par l'objectif}

\subsection{Simulateur}

Le simulateur (environnement gymnasium) représente une cellule de 106 PRB (40~MHz,
espacement 30~kHz) partagée entre eMBB, URLLC et best effort, avec un pas d'une seconde. Sa
capacité, 75~Mb/s, a été mesurée sur la plateforme OAI. Chaque slice est une file fluide : le
trafic arrive, la part de PRB fixe le débit de service, le reste attend dans un tampon. La latence
suit un modèle M/M/1 plus le temps de vidage du tampon, de sorte qu'elle monte avant la
saturation et donne un signal à la politique. Le trafic suit des profils (stable, pic, rampe) avec
du bruit.

L'objectif RL est un schéma pydantic : un poids par slice, un débit minimal par slice, une borne
de latence pour URLLC, et une fenêtre de validité. Il fait partie de l'observation : c'est ce qui
rend la politique \emph{conditionnée par l'objectif}. Le LLM peut changer l'objectif en cours
d'exécution sans ré-entraînement.

\subsection{Entraînement : un premier échec, puis un test valide}

Nous avons d'abord entraîné DDPG, puis TD3 (une variante plus stable de DDPG, prévue comme
solution de repli ; mêmes hyperparamètres dans stable-baselines3). Toutes les politiques battaient
la répartition égale, mais \textbf{le critère de conditionnement fixé à l'avance a échoué} : pour
au moins une graine, la répartition changeait à peine entre deux objectifs. Nous avons gardé cet
échec dans le dossier et cherché les causes, chacune mesurée avant d'être corrigée
(tableau~\ref{tab:causes}).

\begin{table}[!htb]
\centering\small
\begin{tabularx}{\textwidth}{@{}P{5.3cm}L@{}}
\toprule
Cause trouvée & Correction\\
\midrule
Le test lui-même ne séparait pas une politique parfaite d'une mauvaise (cas d'égalité entre slices) & Nouveau test fondé sur la récompense et sur un oracle\\
Déficit réel sous contraintes : sous « protéger URLLC », l'optimum gagne 0,252 par pas, la politique $-0{,}015$ & Ré-étiquetage des objectifs (idée de \emph{hindsight experience replay}) : 80~\% de chaque lot est rejoué sous un autre objectif\\
Une « falaise » dans la récompense : au-delà de deux fois la borne, toute latence coûtait pareil & Pénalité logarithmique $\log(1 + \text{excès}/\text{borne})$\\
Le pic du test ($\times 4$) n'était jamais vu à l'entraînement ($\times 1{,}5$ à 3) & Pics jusqu'à $\times 5$, charge $\times 1{,}2$\\
Le filtre de lissage n'était appliqué qu'au déploiement & Filtre inclus dans l'environnement d'entraînement\\
\bottomrule
\end{tabularx}
\caption{Causes de l'échec du premier critère et corrections.}
\label{tab:causes}
\end{table}

Le nouveau test repose sur une idée simple : \emph{dire à la politique l'objectif A doit l'aider à
bien faire sous A}. On compare la récompense sous A quand on lui annonce A, et quand on lui
annonce un autre objectif B. Un oracle, réservé à l'évaluation, sert de référence : chaque seconde,
il essaie 861 répartitions, prédit la seconde suivante avec les vraies charges futures et applique
la meilleure. Un contrôleur qui ignore l'objectif échoue par construction. Les critères
(avantage positif dans au moins 90~\% des cas, au moins la moitié de l'avantage de l'oracle
capturé, et mieux que la répartition égale) doivent être tenus \emph{par chaque graine}.

Le protocole a été fixé avant tout entraînement : choix de la recette sur un jeu de validation,
confirmation sur cinq graines nouvelles, puis test réservé, joué une seule fois.
\textbf{Les cinq graines passent} (avantage positif dans 97~\% des cas, 70 à 88~\% de l'avantage
de l'oracle, récompense 0,279 à 0,286 par pas contre 0,183 pour la répartition égale). La
figure~\ref{fig:objectifs} montre l'effet : la même politique, sur le même pic URLLC, agit
différemment selon l'objectif annoncé.

\begin{figure}[!htb]
\centering
\includegraphics[width=0.62\textwidth]{learn/img/M02-objectives.png}
\caption{Une politique validée, même pic URLLC (zone grise), deux objectifs. À gauche,
« protéger URLLC » : URLLC prend l'essentiel des PRB en 3~s environ et repasse sous sa borne de
10~ms ; best effort paie. À droite, « maximiser eMBB » : URLLC reçoit peu. Lignes : débit,
latence (échelle log), parts de PRB, récompense.}
\label{fig:objectifs}
\end{figure}

Une seconde série d'entraînement a ensuite ajouté un plancher de 5~\% de PRB par slice, pour
que best effort ne soit jamais affamé ; elle passe aussi le test réservé sur chaque graine. La
politique déployée (« v7 ») vient de cette série.

\subsection{Suivi et déploiement des modèles}

Chaque entraînement et chaque évaluation sont enregistrés dans MLflow (local), avec les
paramètres, le commit git et les courbes. Le modèle déployé est désigné dans un registre par
l'alias \tool{champion}, que seule une règle déterministe peut déplacer : test réservé réussi,
recette réussie sur chaque graine, puis meilleur score de validation. La xApp RL charge ce modèle
une seule fois au démarrage et n'appelle jamais MLflow dans la boucle de contrôle. Les tests
tournent automatiquement à chaque \emph{push} (GitHub Actions).

%----------------------------------------------------------------- 4
\section{Vérificateur déterministe}

Le vérificateur est la règle fixe entre « le LLM l'a proposé » et « le réseau l'applique ». Il
prend l'objectif proposé et l'état courant (horloge, capacité, intentions actives), exécute huit
contrôles dans un ordre fixe, et renvoie \emph{accepté} ou \emph{rejeté} avec une phrase qui
contient les chiffres. Il est écrit en Python pur, sans LLM, sans aléa, sans réseau, sans lecture
de l'horloge (elle lui est passée) : la même entrée donne toujours la même réponse, et un test le
vérifie.

\begin{table}[!htb]
\centering\small
\begin{tabularx}{\textwidth}{@{}P{0.5cm}P{3.6cm}L@{}}
\toprule
& Contrôle & Rejette\\
\midrule
1 & Schéma & Champ inconnu, poids hors $[0,1]$ ou de somme différente de 1, valeur infinie ou NaN\\
2 & Fenêtre & Intention sans fin ; durée restante inférieure à 10~s ou supérieure à un jour\\
3 & Physique de la latence & Borne URLLC $\leq 2$~ms : un slice URLLC vide prend déjà 2~ms\\
4 & Budget & Cibles dont la somme des parts nécessaires dépasse la cellule\\
5 & Plancher & Cibles qui laissent un slice sous 5~\% des PRB\\
6 & Domaine d'entraînement & Cible au-delà de 40~\% de la capacité, borne hors de 5--50~ms (la politique n'a jamais vu ces cas)\\
7 & Conflit & Budget ou plancher violé par la fusion avec les intentions actives ; la plus récente est rejetée\\
8 & Amplitude du changement & Poids qui bouge de plus de 0,5 ou cible de plus de 30~\% de la capacité en une fois\\
\bottomrule
\end{tabularx}
\caption{Les huit contrôles, dans leur ordre d'exécution. Le premier échec donne le motif.}
\end{table}

Pour le budget, chaque cible est convertie en part de PRB. Pour URLLC avec une borne $B$, rester
sous $B$ impose une charge $\rho \leq 1 - \text{base}/B$, donc une cible $T$ demande
$T / (C\,(1 - \text{base}/B))$ de la cellule : 30~Mb/s sous 3~ms demanderaient 120~\% de la
cellule, et sont refusés. Quand plusieurs intentions sont actives, elles sont fusionnées (plus haut
débit minimal par slice, borne la plus stricte, poids de la plus récente), et c'est la
\emph{fusion} qui est vérifiée : ce qui est vérifié est exactement ce que suit la xApp RL. La
figure~\ref{fig:verif} montre la décision pour toutes les paires de cibles.

\begin{figure}[!htb]
\centering
\includegraphics[width=\textwidth]{learn/img/M03-feasibility.png}
\caption{Décision du vérificateur pour chaque paire de cibles (URLLC en abscisse, eMBB en
ordonnée). (a) sans borne de latence ; (b) avec une borne URLLC de 5~ms, qui réduit la zone
acceptée ; (c) avec en plus une intention active « best effort $\geq 20$~Mb/s », qui crée une zone
de conflit. Les couleurs indiquent le contrôle qui rejette.}
\label{fig:verif}
\end{figure}

Les tests couvrent 13 cas acceptés et 24 cas rejetés, chacun avec le contrôle et le motif
attendus, y compris les cas limites (cibles exactement au budget, fenêtres qui se touchent,
arrondis flottants).

%----------------------------------------------------------------- 5
\section{Service LangGraph et intégration}

\subsection{Processus et base partagée}

Les composants tournent comme des processus séparés qui ne communiquent que par une base
SQLite (\tool{var/ric.db}) et un journal d'événements : la xApp RL, le service LangGraph (API
FastAPI et un \emph{worker}), la console de confirmation, et un exportateur Prometheus pour un
tableau de bord Grafana. \tool{POST /intents} ne fait qu'enregistrer la demande ; le
\emph{worker} la traite ensuite. Chaque seconde, la xApp RL lit le segment de politique A1 en
vigueur, construit l'observation (avec la même fonction qu'à l'entraînement), applique la
répartition et écrit les KPI. Si la base est occupée, elle garde le dernier segment et ne
l'attend pas. Une intention qui expire n'a besoin d'aucun processus : l'horloge passe simplement
au segment suivant du calendrier.

\subsection{Boucle de l'agent}

La boucle prévue dans la proposition a été réalisée sous forme d'un graphe LangGraph explicite
(figure~\ref{fig:boucle}). Le LLM extrait d'abord les \emph{faits} de la phrase (slice prioritaire,
débits, borne, durée), puis choisit un outil à la fois. Les chiffres de conformité sont calculés
par du code ; le LLM ne fait que décider « terminé » ou « réviser ».

\begin{figure}[!htb]
\centering
\begin{tikzpicture}[node distance=7mm and 7mm, every node/.style={font=\small}]
  \node[cbox, minimum width=2.4cm] (p) {Analyse\\\footnotesize LLM, schéma JSON};
  \node[cbox, minimum width=2.4cm, right=of p] (pl) {Planification\\\footnotesize un outil à la fois};
  \node[cbox, minimum width=2.4cm, right=of pl] (a) {Action\\\footnotesize \tool{set\_rl\_objective}};
  \node[cbox, minimum width=2.4cm, right=of a, fill=gray!15] (v) {Vérificateur\\\footnotesize 8 contrôles};
  \node[cbox, minimum width=2.4cm, below=14mm of v, dashed] (cf) {Confirmation\\\footnotesize console, démo};
  \node[cbox, minimum width=2.4cm, below=14mm of a] (ap) {Application\\\footnotesize \tool{a1\_policy}};
  \node[cbox, minimum width=2.4cm, below=14mm of pl] (o) {Observation\\\footnotesize 20 s, conformité};
  \node[cbox, minimum width=2.4cm, below=14mm of p] (r) {Réflexion\\\footnotesize LLM : terminé ?};
  \node[cbox, minimum width=2.4cm, below=8mm of r, fill=gray!15] (f) {Fin};
  \node[lbl, above=10mm of p, align=center] (in) {intention reçue\\(\tool{POST /intents})};
  \draw[arr] (in) -- (p);
  \draw[arr] (p) -- (pl); \draw[arr] (pl) -- (a); \draw[arr] (a) -- (v);
  \draw[arr] (v) -- node[lbl,right]{accepté} (cf);
  \draw[arr] (cf) -- (ap);
  \draw[arr] (ap) -- (o); \draw[arr] (o) -- (r);
  \draw[arr] (r) -- node[lbl,right]{oui} (f);
  \draw[arr] (r.north east) -- node[lbl,right,pos=0.35]{réviser (1 fois)} (pl.south west);
  \draw[arr] (v.north) -- ++(0,10mm) -| node[lbl,above,pos=0.25]{rejet avec motif : replanifier (3 rejets au plus)} ([xshift=-5mm]pl.north);
  \draw[arr] ([xshift=-5mm]a.north) to[bend right=50] node[lbl,above]{\tool{get\_kpis}} ([xshift=5mm]pl.north);
\end{tikzpicture}
\caption{Boucle du service LangGraph, telle qu'elle est implémentée. En évaluation automatique,
l'étape de confirmation est sautée.}
\label{fig:boucle}
\end{figure}

Quelques choix de conception :
\begin{itemize}
  \item \textbf{Faits et décisions séparés.} L'analyse ne fait qu'extraire ce qui a été dit ;
        la planification décide. L'extraction peut ainsi être évaluée seule contre des annotations.
  \item \textbf{Une importance de 0 à 10 par slice plutôt que des poids.} Un modèle de 8~milliards
        de paramètres gère mal trois décimales de somme 1 ; la normalisation est faite par du code,
        et le vérificateur contrôle le résultat.
  \item \textbf{En démonstration, les chiffres de l'opérateur ne sont pas modifiés.} Si le LLM
        propose une cible différente de celle demandée, la proposition est renvoyée ; si le
        vérificateur rejette les valeurs demandées, l'intention se termine avec son motif.
  \item \textbf{La confirmation se fait dans une console séparée}, avec un délai d'expiration de
        120~s, et non par un outil du LLM, qui pourrait sinon se l'accorder lui-même.
\end{itemize}

\subsection{Choix du modèle}

Le LLM tourne en local avec Ollama. La règle de choix a été écrite avant les essais : le modèle
doit tenir entièrement dans la mémoire du GPU (8~Go), puis il est classé par le nombre
d'intentions entièrement bien lues sur 20 intentions annotées de réglage, et départagé par la
latence.

\begin{table}[!htb]
\centering\small
\begin{tabularx}{\textwidth}{@{}P{3.2cm}P{2.4cm}P{2.2cm}P{2.6cm}L@{}}
\toprule
Modèle & Intentions correctes & Latence médiane & Part sur le GPU & Décision\\
\midrule
\textbf{\tool{qwen3:8b}} & \textbf{19/20} & \textbf{1,3 s} & \textbf{100~\%} & \textbf{retenu}\\
\tool{qwen3.5:9b} & 19/20 & 2,1 s & 87~\% & exclu (ne tient pas)\\
\tool{ministral-3:8b} & 19/20 & 1,4 s & 89~\% & exclu (ne tient pas)\\
\tool{llama3.1:8b} & 17/20 & 0,9 s & 100~\% & 2 de moins\\
\tool{granite4.2:8b} & 6/20 & 0,4 s & 100~\% & champs numériques omis\\
\bottomrule
\end{tabularx}
\caption{Comparaison des modèles locaux sur le jeu de réglage.}
\end{table}

Les erreurs avaient toutes la même cause (une réponse « compris » sans aucun champ rempli). Une
règle ajoutée au prompt et une vérification dans le code les corrigent : 20/20 sur le jeu de
réglage, un chiffre optimiste par construction puisque le prompt a été ajusté sur ces phrases. Le
score sur des intentions jamais vues sera donné avec l'évaluation de l'interface.

%----------------------------------------------------------------- 6
\section{Plateforme OAI + FlexRIC}

\subsection{Trois slices de bout en bout}

La plateforme tourne dans Docker sous WSL2, sans matériel radio : c\oe ur OAI CN5G v2.2.1,
gNB OAI v3.0.0 en mode rfsim avec son agent E2, Near-RT RIC FlexRIC v3.0.0, et trois nrUE, un par
slice. Les trois slices sont déclarés de façon cohérente dans toutes les couches (AMF, SMF, UPF,
base des abonnés, gNB, UE) ; un test vérifie que les fichiers de configuration concordent.

\begin{table}[!htb]
\centering\small
\begin{tabularx}{\textwidth}{@{}P{2.6cm}P{3.2cm}P{2.2cm}P{2.8cm}L@{}}
\toprule
Slice & S-NSSAI (SST, SD) & DNN & Sous-réseau UE & 5QI\\
\midrule
eMBB & 1, 000001 & \tool{embb} & 10.0.1.0/24 & 9\\
URLLC & 2, 000001 & \tool{urllc} & 10.0.2.0/24 & 7\\
Best effort & 1, 000002 & \tool{be} & 10.0.3.0/24 & 9\\
\bottomrule
\end{tabularx}
\caption{Configuration des slices sur la plateforme.}
\end{table}

Mesures avec iperf3 : environ 71~Mb/s pour un slice seul. L'ensemble consomme environ 4~Go de
RAM et 4 à 7 c\oe urs sur 32, ce qui laisse la place aux xApps, à l'agent et au LLM. Le résultat
principal de cette étape : \textbf{les slices existent dans la signalisation, mais
l'ordonnanceur d'OAI les ignore}. Avec 40~Mb/s offerts à chacun, chaque slice reçoit un tiers de la
cellule et le délai aller-retour monte à environ 16,6~s pour tous, URLLC compris. C'est exactement
le manque que la xApp RL doit combler.

\subsection{Supervision KPM, et un bug du gNB}

La xApp de supervision s'abonne au modèle de service KPM avec un abonnement par slice (filtré
sur le S-NSSAI complet) et trois mesures : débit, PRB utilisés, délai RLC. Une partie en C sur le
RIC imprime les rapports ; une partie en Python les transforme en enregistrements KPI, au même
format que ceux du simulateur. Le service LangGraph et Grafana lisent ainsi les données réelles
sans aucune modification.

Avec un abonnement par slice, OAI v3.0.0 rendait des valeurs absurdes
(figure~\ref{fig:kpm}, à gauche). La cause : le gNB calcule chaque débit comme une différence de
compteurs, et rangeait la « valeur précédente » par position dans la liste d'un abonnement. Avec
un UE par slice, les trois abonnements partageaient la même case, et chaque différence était prise
contre le compteur d'un autre slice. Un correctif d'environ 50 lignes garde la valeur précédente
par UE ; nous construisons notre propre image du gNB. Après correction, KPM suit iperf3 avec une
erreur médiane de 2,5~\%.

\begin{figure}[!htb]
\centering
\includegraphics[width=\textwidth]{learn/img/M07a-stock-vs-patched.png}
\caption{Débit KPM par slice (trait plein) et débit reçu par iperf3 (pointillés), pour
30\,/\,10\,/\,5~Mb/s offerts. À gauche, gNB d'origine (échelle log) : les valeurs n'ont pas de
sens. À droite, gNB corrigé : KPM et iperf3 se superposent.}
\label{fig:kpm}
\end{figure}

Deux remarques sur les mesures : le « délai » mesuré par KPM est le temps d'attente dans la file
RLC du gNB, pas une latence de bout en bout ; et une cellule pleine affiche 74~\% de PRB utilisés,
car la mesure inclut les créneaux montants du motif TDD. Nous corrigeons ce facteur (0,743), qui
se déduit de la configuration du gNB.

\subsection{Contrôle des PRB par E2}

L'étude préalable avait montré que \textbf{le contrôle des PRB par slice n'existe pas dans OAI
v3.0.0} : le modèle de service SLICE de FlexRIC y est une coquille vide. Nous avons donc
modifié le gNB :
\begin{itemize}
  \item un ordonnanceur qui garde chaque slice dans sa propre plage de PRB, en réutilisant
        l'ordonnanceur \emph{proportional fair} d'OAI à l'intérieur de chaque plage ;
  \item un vrai traitement du message de contrôle SLICE, qui refuse les plages hors cellule, qui se
        chevauchent ou de moins de 5~PRB, et renvoie la raison dans l'acquittement.
\end{itemize}
Une petite xApp en C transforme « PRB par slice » en messages de contrôle E2, et
\tool{FlexRICBackend} place la plateforme derrière la même interface que le simulateur. Le gNB
répond en 0,25~ms. Le débit suit les PRB de façon linéaire : 0,654~Mb/s par PRB, de 5 à 96 PRB
(figure~\ref{fig:cap}), soit environ 8~\% de moins que l'hypothèse du simulateur.

\begin{figure}[!htb]
\centering
\includegraphics[width=0.68\textwidth]{learn/img/M07b-cap.png}
\caption{Débit reçu en fonction du nombre de PRB attribués par E2, chaque slice étant surchargé
(80~Mb/s offerts). Trait plein : ajustement mesuré ; tirets : hypothèse du simulateur.}
\label{fig:cap}
\end{figure}

\subsection{La boucle fermée sur la vraie cellule}
\label{sec:boucle}

Nous avons rejoué sur la plateforme le scénario de démonstration du simulateur : un pic URLLC de
7,5 à 30~Mb/s pendant 240~s, soit 90~Mb/s offerts à une cellule d'environ 71~Mb/s. Avec
l'ordonnanceur d'origine, URLLC reçoit 14,7~Mb/s sur 30 ; avec la xApp RL, 28,4~Mb/s, au prix du
best effort, ramené à son plancher de 5~\%.

Enfin, une intention a été soumise pendant ce pic : « Garantis 20~Mb/s au slice URLLC pendant
5 minutes. » Le service LangGraph l'a analysée, le vérificateur l'a acceptée, l'opérateur l'a
confirmée dans la console, et la xApp RL a suivi le nouvel objectif dès la seconde suivante
(figure~\ref{fig:intention}). Une demande précédente à 25~Mb/s avait été rejetée par le
vérificateur : le changement dépassait sa limite d'amplitude (30~\% de la capacité).

\begin{table}[!htb]
\centering\small
\begin{tabularx}{\textwidth}{@{}LP{2.2cm}P{2.2cm}P{2.6cm}P{2.6cm}@{}}
\toprule
Pendant le pic & Part URLLC & Débit URLLC & Secondes à $\geq 20$~Mb/s & File URLLC (médiane)\\
\midrule
Avant l'intention (36 s) & 0,25 & 14,2 Mb/s & 28~\% & 6,3 s\\
Intention en vigueur (184 s) & 0,49 & 31,8 Mb/s & \textbf{95~\%} & \textbf{0,36 s}\\
\bottomrule
\end{tabularx}
\caption{Effet de l'intention sur la vraie cellule (mesures KPM, une exécution).}
\end{table}

\begin{figure}[!htb]
\centering
\includegraphics[width=\textwidth]{learn/img/M07b-intent.png}
\caption{L'intention appliquée sur le gNB réel pendant le pic URLLC (barre grise). Zone teintée :
objectif de l'intention en vigueur. De haut en bas : parts de PRB, débit KPM (pointillé : cible
de 20~Mb/s), délai de file RLC (échelle log). L'oscillation des parts est discutée en
section~\ref{sec:limites}.}
\label{fig:intention}
\end{figure}

%----------------------------------------------------------------- 7
\clearpage
\section{Difficultés et limites}
\label{sec:limites}

\begin{table}[H]
\centering\small
\begin{tabularx}{\textwidth}{@{}P{4.2cm}LL@{}}
\toprule
Difficulté & Ce que nous avons fait & Ce qui reste\\
\midrule
Le premier critère de conditionnement par l'objectif a échoué & Causes mesurées (test mal conçu, quatre défauts d'entraînement), nouveau test fixé avant l'entraînement, réussi sur chaque graine & Validé en simulation uniquement\\
Bug KPM d'OAI v3.0.0 avec plusieurs abonnements & Correctif d'environ 50 lignes, image du gNB maison & Correctif non proposé à OAI\\
Pas de contrôle des PRB par E2 dans OAI v3.0.0 & Ordonnanceur par slice et traitement du contrôle SLICE ajoutés au gNB & gNB à reconstruire depuis les sources ; les PRB réservés et inutilisés sont perdus (best effort seul passe de 64,7 à 19,7~Mb/s à côté d'un eMBB inactif)\\
FlexRIC ne transmet pas la réponse du gNB à la xApp & Le refus reste visible dans le journal du gNB et dans les captures & La xApp ne distingue pas un ordre accepté d'un ordre refusé\\
La politique oscille sur la plateforme (part eMBB entre 0,22 et 0,84, période 10--12~s) & Cause trouvée : la politique voit l'effet de ses décisions en retard (une seconde mesurée, deux d'après la simulation), ce qu'elle n'a jamais vu à l'entraînement. Lissage plus fort : part de la cellule déplacée par seconde de 10,8 à 4,3~\%, secondes avec URLLC $\geq 20$~Mb/s de 69 à 90~\% & L'oscillation diminue sans disparaître ; le délai de file URLLC reste de l'ordre de la seconde. Un ré-entraînement avec retard est à décider\\
KPM ne mesure pas la charge offerte & Sur la plateforme, la xApp RL lit le débit que le script de trafic envoie ; sans cette information, elle se fige après le pic & Estimer la demande à partir de mesures du gNB (tampon RLC)\\
Deux exécutions de la même politique sur le même scénario se sont comportées différemment & Constaté et noté & Répéter chaque mesure (prévu en phase 5)\\
Plateforme simplifiée & Canal rfsim idéal, un UE par slice & Résultats à lire comme une validation de la chaîne, pas comme une mesure radio réaliste\\
\bottomrule
\end{tabularx}
\caption{Difficultés rencontrées et limites actuelles.}
\end{table}

%----------------------------------------------------------------- 8
\section{Suite du travail}

Les deux phases suivantes reprennent le plan de la proposition et seront présentées dans le
prochain rapport :
\begin{itemize}
  \item \textbf{Phase 4, interface :} passerelle MCP qui expose exactement
        \tool{submit\_intent}, \tool{get\_kpis} et \tool{get\_status} ; Hermes Agent isolé dans
        un conteneur, outils généralistes désactivés ; tests de demandes hors périmètre.
  \item \textbf{Phase 5, résultats :} comparaison des quatre contrôleurs (répartition statique,
        RL seul, LLM seul, LLM + RL) sur les mêmes scénarios, en simulation et sur la plateforme ;
        évaluation de l'interface (extraction sur des intentions annotées jamais vues, appels
        d'outils, tentatives hors périmètre).
\end{itemize}
Les limites de la section~\ref{sec:limites}, en particulier l'oscillation et la charge offerte,
seront traitées en parallèle.

%----------------------------------------------------------------- refs
\section*{Références}
\small
\begin{enumerate}[label={[\arabic*]}]
  \item Rezazadeh et al., apprentissage par renforcement profond multi-agents pour le slicing du RAN.
  \item Wu et al., IA agentique pilotée par LLM pour la résilience en O-RAN.
  \item Bao et al., LLM-hRIC : contrôle hiérarchique LLM + RL dans le RIC.
  \item Fujimoto et al., \emph{Addressing Function Approximation Error in Actor-Critic Methods} (TD3), ICML 2018.
  \item Andrychowicz et al., \emph{Hindsight Experience Replay}, NeurIPS 2017.
  \item OpenAirInterface : gNB et nrUE v3.0.0, OAI CN5G v2.2.1 (\url{https://openairinterface.org}).
  \item FlexRIC v3.0.0 (\url{https://gitlab.eurecom.fr/mosaic5g/flexric}).
  \item O-RAN Alliance, E2SM-KPM v3.00 ; 3GPP TS 28.552 (mesures de performance).
  \item LangGraph, Ollama, stable-baselines3, gymnasium, MLflow.
\end{enumerate}

\end{document}
